% Options for packages loaded elsewhere
\PassOptionsToPackage{unicode}{hyperref}
\PassOptionsToPackage{hyphens}{url}
%
\documentclass[
]{book}
\usepackage{amsmath,amssymb}
\usepackage{iftex}
\ifPDFTeX
  \usepackage[T1]{fontenc}
  \usepackage[utf8]{inputenc}
  \usepackage{textcomp} % provide euro and other symbols
\else % if luatex or xetex
  \usepackage{unicode-math} % this also loads fontspec
  \defaultfontfeatures{Scale=MatchLowercase}
  \defaultfontfeatures[\rmfamily]{Ligatures=TeX,Scale=1}
\fi
\usepackage{lmodern}
\ifPDFTeX\else
  % xetex/luatex font selection
\fi
% Use upquote if available, for straight quotes in verbatim environments
\IfFileExists{upquote.sty}{\usepackage{upquote}}{}
\IfFileExists{microtype.sty}{% use microtype if available
  \usepackage[]{microtype}
  \UseMicrotypeSet[protrusion]{basicmath} % disable protrusion for tt fonts
}{}
\makeatletter
\@ifundefined{KOMAClassName}{% if non-KOMA class
  \IfFileExists{parskip.sty}{%
    \usepackage{parskip}
  }{% else
    \setlength{\parindent}{0pt}
    \setlength{\parskip}{6pt plus 2pt minus 1pt}}
}{% if KOMA class
  \KOMAoptions{parskip=half}}
\makeatother
\usepackage{xcolor}
\usepackage{longtable,booktabs,array}
\usepackage{calc} % for calculating minipage widths
% Correct order of tables after \paragraph or \subparagraph
\usepackage{etoolbox}
\makeatletter
\patchcmd\longtable{\par}{\if@noskipsec\mbox{}\fi\par}{}{}
\makeatother
% Allow footnotes in longtable head/foot
\IfFileExists{footnotehyper.sty}{\usepackage{footnotehyper}}{\usepackage{footnote}}
\makesavenoteenv{longtable}
\usepackage{graphicx}
\makeatletter
\def\maxwidth{\ifdim\Gin@nat@width>\linewidth\linewidth\else\Gin@nat@width\fi}
\def\maxheight{\ifdim\Gin@nat@height>\textheight\textheight\else\Gin@nat@height\fi}
\makeatother
% Scale images if necessary, so that they will not overflow the page
% margins by default, and it is still possible to overwrite the defaults
% using explicit options in \includegraphics[width, height, ...]{}
\setkeys{Gin}{width=\maxwidth,height=\maxheight,keepaspectratio}
% Set default figure placement to htbp
\makeatletter
\def\fps@figure{htbp}
\makeatother
\setlength{\emergencystretch}{3em} % prevent overfull lines
\providecommand{\tightlist}{%
  \setlength{\itemsep}{0pt}\setlength{\parskip}{0pt}}
\setcounter{secnumdepth}{5}
\usepackage{booktabs}
\ifLuaTeX
  \usepackage{selnolig}  % disable illegal ligatures
\fi
\usepackage[]{natbib}
\bibliographystyle{plainnat}
\IfFileExists{bookmark.sty}{\usepackage{bookmark}}{\usepackage{hyperref}}
\IfFileExists{xurl.sty}{\usepackage{xurl}}{} % add URL line breaks if available
\urlstyle{same}
\hypersetup{
  pdftitle={Apostila de Introdução ao R},
  pdfauthor={Profª. Angélica Maria T. Ribeiro},
  hidelinks,
  pdfcreator={LaTeX via pandoc}}

\title{Apostila de Introdução ao R}
\author{Profª. Angélica Maria T. Ribeiro}
\date{2024-01-22}

\begin{document}
\maketitle

{
\setcounter{tocdepth}{1}
\tableofcontents
}
\hypertarget{introduuxe7uxe3o}{%
\chapter{Introdução}\label{introduuxe7uxe3o}}

\hypertarget{sobre-o-r-e-rstudio}{%
\section{Sobre o R e RStudio}\label{sobre-o-r-e-rstudio}}

\hypertarget{instalauxe7uxe3o}{%
\section{Instalação}\label{instalauxe7uxe3o}}

\hypertarget{estrutura-buxe1sica-do-r-e-rstudio}{%
\section{Estrutura básica do R e RStudio}\label{estrutura-buxe1sica-do-r-e-rstudio}}

\hypertarget{introduuxe7uxe3o-uxe0-linguagem-r}{%
\section{Introdução à linguagem R}\label{introduuxe7uxe3o-uxe0-linguagem-r}}

\hypertarget{estruturas-de-dados}{%
\chapter{Estruturas de Dados}\label{estruturas-de-dados}}

\hypertarget{vetores-e-fatores}{%
\section{Vetores e fatores}\label{vetores-e-fatores}}

\hypertarget{matrizes-e-data-frames.}{%
\section{Matrizes e data frames.}\label{matrizes-e-data-frames.}}

\hypertarget{operauxe7uxf5es-com-vetores-e-matrizes}{%
\section{Operações com vetores e matrizes}\label{operauxe7uxf5es-com-vetores-e-matrizes}}

\hypertarget{listas}{%
\section{Listas}\label{listas}}

\hypertarget{importauxe7uxe3o-e-manipulauxe7uxe3o-de-dados}{%
\section{Importação e Manipulação de dados}\label{importauxe7uxe3o-e-manipulauxe7uxe3o-de-dados}}

  \bibliography{book.bib}

\end{document}
